% Page 032

\seite{32}

\margmark{Jesus! Maria! Jo}%
seph! Valentin!         Zum frommen Andenken

\pend
\begin{verse}
an unseren innigstgeliebten Sohn,\\
Bruder, Neffen und Vetter\\
den Theologen\\
Valentin Stoffels
\end{verse}
\pstart
Leutnant der Reserve im Inf.-Rgt. 97,\ob
Inhaber des Eisernen Kreuzes II. Klasse.

Er war geboren am 24.\ September 1891\ob
in Sefferweich, trat nach Absolvierung des\ob
Gymnasiums in Prüm zu Ostern in das\ob
Bischöfliche Priesterseminar ein und wid\ohb
mete sich mit heißer Liebe dem hin und\unsicher\ob
guten Erfolge der Theologie, begann 1914\ob
bruch des Krieges. Bis zum 18.\ Mai 1915 war er in einem nach\ohb
garten Stelle [...] der Winterstellung in Rußland mit.\ob
Im Felde angenommen der\ob
Krankheit wieder hergestellt, kam er im\ob
Herbst 1917 als Leutnant wieder ins Feld\ob
und fiel an der Spitze seines Zuges von\ob
tödlichen dem Bataillon schwer verwundet\ob
in die Hände.

Das feierliche Levitenamt findet statt\ob
am 14.\ Oktober, vormittags 10 Uhr in der\ob
Kirche zu Sefferweich.\ob
Sefferweich und im Felde, 10.\ Okt. 1918.\ob
Die tieftrauernden Eltern und Geschwister.
\pend
